\section{Verification evidence, reproduction and scope}\label{app:verification}
\subsection{What was checked, and when}
The formal proof is the checked joint project of 29 October 2026. The repository's root directory \texttt{lean/} contains the
complete project: all 326 source modules, \texttt{lean-toolchain},
\texttt{lakefile.toml} and \texttt{lake-manifest.json}. Every project
source hash matches the earlier successful independent source build.
The new statement map was made by inspecting the definitions and
actual theorem declarations, and comparing them with the preserved
compiler-printed final types, rather than relying on comments alone.

The earlier source build compiled all 326 project modules independently,
including the 319-module closure of \path{CDCert.Main} and
\path{CDCert.TwoInputs}. The earlier 20-result axiom inspection passed
with only \texttt{propext}, \texttt{Classical.choice} and
\texttt{Quot.sound}. These are standard Lean foundations. There is no
\texttt{sorry}, project-specific axiom or \texttt{native\_decide} in
the checked source project. In particular the final theorem does not
assume the finite inequality, the free-to-locked injection, or the
correctness of a native program.

The fresh correspondence check on 1 November 2026 succeeded.
\path{audit/TutorialCorrespondence.lean} checks 23 explicit
restatements: the three semantic interfaces and all twenty audited
results. \path{audit/TutorialSupportStatements.lean} checks 31 further
restatements, covering completion, deletion, fibres, moves, components,
concatenation, the base values and the exact rational form of the finite
inequality. All 54 compile without added assumptions.
\path{audit/TutorialNamedInterfaces.lean} checks all 190 original
definitions and helper declarations named in the proof, prints the
bodies of the 72 definitions, and checks every axiom list. Only the
standard foundations listed above occur.

These correspondence checks were compiled against hash-verified proof
objects from the earlier complete source build. Their successful logs,
wrapper sources and axiom lists are supplied with the tutorial. This
expository revision preserves the original formal proof and does not
claim another complete Lean source build.

The files \path{audit/SOURCE_IDENTITY.json},
\path{audit/TUTORIAL_LEAN_MAP.json},
\path{audit/ENTRYPOINT_SIGNATURES.txt} and
\path{audit/prior_kernel_build/AXIOM_CHECK_RECEIPT.json}
record the original evidence.
\path{audit/FORMAL_STATEMENT_COVERAGE.json} maps every new restatement
to its tutorial statement or definition. The successful current logs,
source hashes and axiom lists are recorded in
\path{audit/fresh_check/CORRESPONDENCE_VERIFICATION_RECEIPT.json}.
The corresponding human proofs, finite
normalisations and named helper interfaces are given beside each
statement in this tutorial.

\subsection{Rebuilding the formal proof}
Use Lean 4.26.0 and Mathlib v4.26.0 with the transitive revisions pinned
in the manifest. From the repository root:
\begin{quote}\small\ttfamily
cd lean\\
lake exe cache get\\
lake build CDCert.Main CDCert.TwoInputs\\
lake env lean receipts/FinalAxioms.lean
\end{quote}
The cache command obtains compiled standard libraries. The build
command then checks the project sources. Its large finite
kernel computations may take substantial time and memory. A successful
result requires completion, not a partially built directory or an
import that timed out. To check the additional correspondence source
once that build is available, run from the same \texttt{lean/} directory:
\begin{quote}\small\ttfamily
lake env lean ../tutorial/audit/TutorialCorrespondence.lean\\
lake env lean ../tutorial/audit/TutorialSupportStatements.lean\\
lake env lean ../tutorial/audit/TutorialNamedInterfaces.lean
\end{quote}
The complete formal proof is independent of executing the native
Python/C++ generators. The repository also contains the native generators, exact data and
separately dated replay/regeneration records.
They are not needed as trust assumptions of the Lean theorem.

\subsection{Written proof versus a literal formal statement}
A formal proof certifies a precise statement relative to its definitions.
Here the semantic definitions quantify over actual strict total orders
and every odd voter profile. The order/list equivalence, correspondence
of laws and fixed pairs, and preservation of cardinality and swaps are
proved inside the project. The tutorial's real and rational notation is
matched to the natural/integer statements by the explicit index and
scaling tables. These are the substantive safeguards against proving a
bound for a different class or a differently normalised statistic.

A human proof may use equivalent steps: the prose amplification uses
arbitrary \(t\) and a limit; Lean uses powers of two and an integer
contradiction. The hypotheses and conclusion agree. The correspondence
paragraphs explain such differences.

The following remain outside the literal theorem coverage: intuitive
commentary, exact labelled examples checked by a separate small script,
search history and performance, bibliographical claims, attribution and
novelty. The twelve-idea division is an organisational choice, not a
formal theorem. No priority claim follows from successful formalisation.
The numerical comparison tables are exact integer calculations from
the proved recurrences, with long displayed entries rounded upwards.
The optional classified starting values are explicitly external
inputs, and no additional Lean theorem is claimed for them or for
the individual table entries. The new examples and tables are checked
separately by \texttt{scripts/build\_numerical\_examples.py}.

\subsection{Sources and credit}
The mathematical source is the joint paper by P. Moros and S. Riis,
\emph{Improved Upper Bounds on the Size of Condorcet Domains}~\cite{joint},
now available as the review draft of 1 November 2026 in the accompanying
repository. Its earlier source and the unchanged formal proof are preserved
in the checked snapshot of 29 October. Moros's earlier note~\cite{moros}
is an AI-assisted precursor within the same investigation. The joint paper
presents the improved \(32/11\) bound and its component reduction,
together with the peak-pit bound \(8/3\).

This tutorial is an exposition by S\o ren Riis of that joint work.
Claude (Anthropic) and ChatGPT/Codex (OpenAI) assisted the exploration,
checking and drafting of both the research and explanatory material.
Computational optimisation selected finite correction coefficients;
their validity is established by exact verification. The tutorial's
byline identifies responsibility for this exposition, while the
underlying results and formalisation remain credited to their sources.
The sources also credit Fathi's counting observations and the
majority/median background of Puppe and Slinko.

The following background references orient the general mathematical
ideas; the proof above and its explicit certificate provide the actual
logical ingredients for this tutorial.
\clearpage
\begin{thebibliography}{9}\small\raggedright
\bibitem{joint} P. Moros and S. Riis, \emph{Improved Upper Bounds on the
Size of Condorcet Domains}, joint review draft, 1 November 2026;
complete accompanying Lean formalisation and computational material.
\url{https://github.com/SR123/condorcet-upper-bound}.
The formal proof is pinned to snapshot \texttt{0d5dd168854b}
(29 October 2026).
\bibitem{moros} P. Moros, \emph{Condorcet domains have fewer than
\(2.95^n\) orders}, earlier AI-assisted research note (2026).
\url{https://doi.org/10.5281/zenodo.23654773}.
\bibitem{sen} A. K. Sen, A possibility theorem on majority decisions,
\emph{Econometrica} 34 (1966), 491--499.
\bibitem{median} C. Puppe and A. Slinko, Condorcet domains, median graphs
and the single-crossing property, \emph{Economic Theory} 67 (2019),
285--318. \url{https://doi.org/10.1007/s00199-017-1084-6}.
\bibitem{discharging} D. W. Cranston and D. B. West, An introduction to
the discharging method via graph coloring, \emph{Discrete Mathematics}
340 (2017), 766--793. \url{https://arxiv.org/abs/1306.4434}.
\bibitem{certificates} K. K. H. Cheung, A. Gleixner and D. E. Steffy,
Verifying integer programming results, \emph{IPCO 2017}, 148--160.
\url{https://arxiv.org/abs/1611.08832}.
\bibitem{smallcd} D. Akello-Egwell, C. Leedham-Green, A. Litterick,
K. Markstr\"om and S. Riis, \emph{Condorcet Domains of Degree at most
Seven}, arXiv:2306.15993, version 5 (2023), Section 4.1 and Tables 1--4.
\url{https://arxiv.org/abs/2306.15993}.
\bibitem{eightcd} C. Leedham-Green, K. Markstr\"om and S. Riis,
The largest Condorcet domain on 8 alternatives,
\emph{Social Choice and Welfare} 62 (2024), 109--116, Theorem 4.1.
\url{https://doi.org/10.1007/s00355-023-01481-3}.
\end{thebibliography}
